\title{Group Sequence Policy Optimization}

\begin{abstract}
We present Group Sequence Policy Optimization (GSPO), a reinforcement learning framework specifically designed for robust, effective, and efficient training of large language models. Departing from prior methods that rely on token-level importance ratios, GSPO calculates the importance ratio over entire sequences, applying sequence-level mechanisms for clipping, reward assignment, and optimization. Our results demonstrate that GSPO outperforms the GRPO algorithm in both training effectiveness and overall model performance, delivers considerable stability for Mixture-of-Experts (MoE) RL training, and can streamline RL system architecture. These advantages of GSPO have directly contributed to the significant advancements seen in the most recent Qwen3 models.
\end{abstract}

\section{Introduction}
\label{sec:intro}

Reinforcement learning (RL) has become a central approach for advancing the capabilities of language models \citep{o1,dpsk-r1,qwq32b,qwen3}. By leveraging RL at scale, these models acquire the skills required to address complex tasks—including competition-level mathematics and programming—by engaging in more elaborate and extended reasoning.

For scaling RL effectively through increased computational resources, it is critical to ensure that training procedures remain both stable and resilient. Yet, leading RL algorithms such as GRPO \citep{grpo} frequently encounter pronounced stability challenges when applied to very large language models, which can precipitate catastrophic, irrecoverable collapses in model performance \citep{qwen3, minimax-m1}. Such instability significantly obstructs attempts to further advance the capabilities of language models via extended RL training.

This work demonstrates that the root cause of instability in GRPO lies in the improper use and subsequent breakdown of importance sampling weights within its algorithmic structure. This flaw results in high-variance noise during training, which builds up as response length increases and is exacerbated by the application of the clipping strategy, eventually leading to the collapse of the model.

In response to these fundamental challenges, we introduce \textbf{Group Sequence Policy Optimization (GSPO)}, an RL framework tailored for the development of large language models. The principal novelty of GSPO is its rigorously derived formulation of the importance ratio, which utilizes sequence likelihood in line with the essential concept of importance sampling \citep{click}. Furthermore, GSPO leverages normalized rewards—interpreted as the advantages—across several responses per query, thereby promoting consistency between reward assignment at the sequence level and policy optimization.

Through our empirical analysis, we establish that GSPO markedly surpasses GRPO in terms of training robustness, computational efficiency, and achieved results. Notably, GSPO naturally addresses the stability difficulties commonly encountered during RL training with large-scale Mixture-of-Experts (MoE) architectures, obviating reliance on intricate stabilization methods and thereby indicating prospects for streamlining RL system design. These strengths of GSPO have, in turn, played a pivotal role in driving substantial performance gains in the most recent Qwen3 models. We anticipate that GSPO will serve as a dependable and extensible algorithmic basis, fostering sustained progress in large-scale RL training with language models.

\section{Preliminaries}

\paragraph{Notation}
Throughout this work, we refer to an autoregressive language model with parameters $\theta$ as a policy, denoted by $\pi_\theta$.
Let $x$ represent a query, and let $\mathcal{D}$ be the collection of all queries. For a specific response $y$ to a query $x$, the probability assigned by the policy $\pi_\theta$ is represented as $\pi_\theta (y | x)=\prod_{t=1}^{|y|} \pi_\theta (y_t | x, y_{<t} )$, where $|y|$ indicates the total number of tokens in $y$. The quality of a query-response pair $(x, y)$ can be evaluated using a verifier $r$, which assigns a reward $r(x, y) \in [0, 1]$.

\paragraph{Proximal Policy Optimization (PPO)} Drawing on trajectories sampled from the previous policy $\pi_{\theta_\text{old}}$, PPO \citep{ppo} limits the magnitude of policy changes by introducing a clipping strategy that confines updates to a local neighborhood around the old policy. The algorithm maximizes the following policy objective (for simplicity, we exclude the KL regularization term, as it is not central to the present work):

\begin{align}
\mathcal{J}_\text{PPO}(\theta) = \mathbb{E}_{ x \sim \mathcal{D},\, y \sim \pi_{\theta_\text{old}}( \cdot | x) }
\left[ \frac{1}{|y|} \sum_{t=1}^{|y|} 
\min \left( w_{t}(\theta) \widehat{A}_{t},  \, \mathrm{clip} \left( w_{t}(\theta), 1 - {\varepsilon}, 1 + {\varepsilon}\right) \widehat{A}_{t} \right)
\right],
\end{align}

Here, the token importance ratio $w_{t}(\theta)$ is given by $
w_{t}(\theta) = \frac{ \pi_{\theta} (y_{t} | x, y_{<t}) }{ \pi_{\theta_\text{old}} (y_{t} | x,y_{<t})} $, with the advantage $\widehat{A}_{t}$ associated with $y_t$ calculated using a separate value model. The parameter $\varepsilon$ specifies the allowed range for clipping the importance ratios.

A primary practical difficulty with PPO is its significant dependence on the value model. Notably, the value model is typically as large as the policy model, resulting in substantial demands on both memory and computational resources. Additionally, the overall success of the algorithm is contingent upon the accuracy of its value predictions. Obtaining a trustworthy value model is already a complex undertaking; scaling such a model to handle lengthier outputs and more intricate problem settings amplifies this complexity.

\paragraph{Group Relative Policy Optimization (GRPO)} GRPO \citep{grpo} eliminates reliance on a value function by assessing the relative advantage of individual responses among a collection of candidate responses for a single query. More precisely, GRPO seeks to maximize the following objective:

\begin{align}
\label{equ:grpo}
\mathcal{J}_\text{GRPO}(\theta) = \mathbb{E}_{ x \sim \mathcal{D},\, \{y_i\}_{i=1}^G \sim \pi_{\theta_\text{old}}( \cdot | x) }
\left[ \frac{1}{G} \sum_{i=1}^{G} \frac{1}{|y_i|} \sum_{t=1}^{|y_i|} 
\min \left( w_{i,t}(\theta) \widehat{A}_{i,t},  \, \mathrm{clip} \left( w_{i,t}(\theta), 1 - {\varepsilon}, 1 + {\varepsilon}\right) \widehat{A}_{i,t} \right)
\right],
\end{align}

where $G$ is the number of generated responses to each query $x$ (i.e., the group size), and the importance ratio $w_{i,t}(\theta)$ and advantage $\widehat{A}_{i,t}$ of token $y_{i,t}$ are:

\begin{align}
    w_{i,t}(\theta) = \frac{ \pi_{\theta}(y_{i,t} \mid x, y_{i,<t}) }{ \pi_{\theta_\text{old}}(y_{i,t} \mid x, y_{i,<t}) },\qquad
    \widehat{A}_{i,t} = \widehat{A}_i = \frac{ r(x, y_i) - \operatorname{mean}\left(\{ r(x, y_i) \}_{i=1}^G\right) }{ \operatorname{std}\left(\{ r(x, y_i) \}_{i=1}^G\right) },
\end{align}

respectively, where all the tokens in $y_i$ share the same advantage as $\widehat{A}_{i}$.

\section{Motivation}
\label{sec:motivation}

As models become larger, incorporate greater sparsity (such as found in Mixture-of-Experts architectures), and generate longer responses, reinforcement learning workloads require larger rollout batch sizes to fully leverage hardware capabilities. To enhance the efficiency with which samples are utilized, it is common to split these extensive rollout batches into several smaller mini-batches for performing gradient updates. However, this strategy results in an off-policy learning scenario: the responses $y$ are typically drawn from an earlier version of the policy, $\pi_{\theta_\text{old}}$, instead of the current policy $\pi_\theta$ that is the focus of optimization. Consequently, this phenomenon also highlights why clipping mechanisms are integral to both PPO and GRPO—to restrict the influence of excessively ``off-policy'' samples during the computation of gradient estimates.

Although approaches such as clipping attempt to address the off-policy mismatch, our analysis reveals a deeper challenge within GRPO: \textit{the objective itself is not well-defined}. This inadequacy is especially pronounced when scaling up to larger models and training them on tasks involving extended responses, often resulting in severe model collapse. The root of this ill-posedness lies in the improper application of importance sampling weights. Importance sampling fundamentally seeks to compute the expectation of a function $f$ under a desired target distribution $\pi_\text{tar}$ by appropriately adjusting the contributions of samples obtained from a behavior distribution $\pi_\text{beh}$:

\begin{align}
\mathbb{E}_{z \sim \pi_\text{tar}} \left[ f(z) \right] = \mathbb{E}_{z \sim \pi_\text{beh}} \left[ \frac{ \pi_\text{tar} (z) }{ \pi_\text{beh} (z)} \, f(z) \right].
\end{align}

Importantly, effective correction of distributional discrepancies using the importance weight $\frac{ \pi_\text{tar} (z) }{ \pi_\text{beh} (z)}$ depends on averaging across a large number of samples ($N \gg 1$) drawn from the behavior distribution $\pi_\text{beh}$.

Conversely, GRPO incorporates the importance weight $\frac{ \pi_{\theta} (y_{i,t} | x, y_{i,<t}) }{ \pi_{\theta_\text{old}} (y_{i,t} | x,y_{i,<t})}$ at each time step $t$ of the token sequence. This weighting is computed from a single sample $y_{i,t}$ drawn from the next-token distribution $\pi_{\theta_\text{old}}( \cdot | x, y_{i,<t})$ for each position, resulting in an inability to properly correct for distribution mismatch as originally intended. Rather than correcting, this strategy injects substantial high-variance noise into the gradient estimates during training. The resulting variance accumulates across lengthy sequences and is worsened further by the presence of clipping. In our experiments, we have found that this effect can precipitate a complete collapse of the model—a failure mode from which recovery is nearly impossible. After such a collapse, further training efforts, such as reverting to earlier checkpoints, adjusting hyperparameters like clipping thresholds, increasing sequence lengths, or altering RL query schemes, have proven ineffective in restoring the model.

The preceding analysis highlights a basic flaw in the construction of GRPO. Specifically, the ineffectiveness of token-level importance weighting reveals a key insight: \textbf{optimization objectives must be defined at the same level as the reward}. Given that rewards are assigned per sequence, implementing off-policy corrections at the token level introduces complications. Therefore, we are prompted to abandon the token-based objective in favor of leveraging importance weights and executing optimization at the \textit{sequence level} instead.

\section{Algorithm}

\subsection{GSPO: Group Sequence Policy Optimization}

Although the use of token-level importance weights $\frac{ \pi_{\theta} (y_{i,t} | x, y_{i,<t}) }{ \pi_{\theta_\text{old}} (y_{i,t} | x,y_{i,<t})}$ presents challenges within GRPO, our analysis reveals that the \textit{sequence-level} importance weight $\frac{ \pi_{\theta} (y | x) }{ \pi_{\theta_\text{old}} (y | x)}$ acquires a more straightforward theoretical interpretation in language generation tasks: it quantifies the extent to which a response $y$ drawn from $\pi_{\theta_\text{old}} (\cdot | x)$ diverges from $\pi_{\theta} (\cdot | x)$. This is consistent with sequence-level rewards and offers an informative signal regarding the application of the clipping mechanism.

Based on this straightforward observation, we propose the \textbf{Group Sequence Policy Optimization (GSPO)} algorithm. GSPO employs the following sequence-level optimization objective:

\begin{align}
\mathcal{J}_\text{GSPO} (\theta) = 
\mathbb{E}_{ x \sim \mathcal{D},\, \{y_i\}_{i=1}^G \sim \pi_{\theta_\text{old}}( \cdot | x) }
\left[ 
\frac{1}{G} \sum_{i=1}^{G}
\min \left( s_{i}(\theta)  \widehat{A}_{i},\, \mathrm{clip}\big( s_{i}(\theta), 1 - {\varepsilon}, 1 + {\varepsilon} \big)\widehat{A}_{i} \right)
\right],
\label{equ:gspo}
\end{align}

where we adopt the group-based advantage estimation:

\begin{align}
\widehat{A}_{i} = \frac{r(x, y_i) - \mathrm{mean} \left( \{ r(x, y_i) \}_{i=1}^G \right) }{ \mathrm{std} \left( \{ r(x, y_i) \}_{i=1}^G \right) },
\end{align}

and define the importance ratio $s_{i}(\theta)$ based on sequence likelihood \citep{click}:

\begin{align}
s_{i}(\theta) = \left( \frac{ \pi_{\theta} (y_i | x) }{ \pi_{\theta_\text{old}} (y_i | x)} \right)^{\frac{1}{|y_i|}}
=
\exp \left( \frac{1}{|y_i|} \sum_{t=1}^{|y_i|} \log \frac{ \pi_{\theta} (y_{i,t} | x, y_{i,<t}) }{ \pi_{\theta_\text{old}} (y_{i,t} | x,y_{i,<t})} \right).
\end{align}

As a result, GSPO implements clipping at the response level rather than the token level, thereby filtering out highly “off-policy” examples from contributing to the gradient estimation. This approach ensures consistency with both sequence-level reward assignment and optimization strategies. Additionally, we employ length normalization in $s_{i}(\theta)$ to both lower its variance and maintain $s_{i}(\theta)$ within a standardized numerical interval. Absent such normalization, substantial variations in the probabilities assigned to a handful of tokens could produce large, erratic shifts in the sequence-wise importance ratio; this would also necessitate different clipping thresholds for responses depending on their lengths. Furthermore, we emphasize that due to varying definitions of importance ratios, the magnitude of clipping ranges utilized in GSPO tends to differ significantly—often by an order of magnitude—from those used in earlier methods such as GRPO.

\subsection{Gradient Analysis}
\label{subsec:gradient}

We can derive the gradient of the GSPO objective as follows (clipping is omitted for brevity):

\begin{align}
\nabla_{\theta} \mathcal{J}_\text{GSPO} (\theta)
=&\ 
\nabla_{\theta} \mathbb{E}_{ x \sim \mathcal{D},\, \{y_i\}_{i=1}^G \sim \pi_{\theta_\text{old}}( \cdot | x) }
\left[ 
\frac{1}{G} \sum_{i=1}^{G} s_{i}(\theta) \widehat{A}_{i}
\right] \\
= &\ 
\mathbb{E}_{ x \sim \mathcal{D},\, \{y_i\}_{i=1}^G \sim \pi_{\theta_\text{old}}( \cdot | x) }
\left[ 
\frac{1}{G} \sum_{i=1}^{G}
s_{i}(\theta) \widehat{A}_{i}
\cdot \nabla_{\theta} \log s_{i}(\theta)
\right] \\
=&\ 
\mathbb{E}_{ x \sim \mathcal{D},\, \{y_i\}_{i=1}^G \sim \pi_{\theta_\text{old}}( \cdot | x) }
\left[ 
\frac{1}{G} \sum_{i=1}^{G} 
\left( \frac{ \pi_{\theta} (y_i | x) }{ \pi_{\theta_\text{old}} (y_i | x)} \right)^{\frac{1}{|y_i|}} \widehat{A}_{i} 
\cdot \frac{1}{|y_i|} \sum_{t=1}^{|y_i|} \nabla_{\theta} \log \pi_{\theta} (y_{i,t} | x, y_{i,<t}) 
\right].
\label{equ:gspo_grad}
\end{align}

For comparison, the gradient of the GRPO objective is as follows (note that $\widehat{A}_{i,t} = \widehat{A}_{i}$):

\begin{align}
\nabla_{\theta} \mathcal{J}_\text{GRPO} (\theta) 
=&\ 
\nabla_{\theta} \mathbb{E}_{ x \sim \mathcal{D},\, \{y_i\}_{i=1}^G \sim \pi_{\theta_\text{old}}( \cdot | x) }
\left[ \frac{1}{G} \sum_{i=1}^{G} \frac{1}{|y_i|} \sum_{t=1}^{|y_i|} 
w_{i,t}(\theta) \widehat{A}_{i,t}
\right] \\
=&\
\mathbb{E}_{ x \sim \mathcal{D},\, \{y_i\}_{i=1}^G \sim \pi_{\theta_\text{old}}( \cdot | x) }
\left[ \frac{1}{G} \sum_{i=1}^{G} \widehat{A}_{i} 
\cdot \frac{1}{|y_i|} \sum_{t=1}^{|y_i|} 
\frac{ \pi_{\theta} (y_{i,t} | x, y_{i,<t}) }{ \pi_{\theta_\text{old}} (y_{i,t} | x,y_{i,<t})} 
\nabla_{\theta} \log \pi_{\theta} (y_{i,t} | x, y_{i,<t})  
\right].
\end{align}

Consequently, the main difference between GSPO and GRPO concerns \textit{the manner in which they apply weights to the gradients of the tokens’ log likelihoods}. GRPO assigns each token a weight corresponding to its so-called ``importance weight'' $\frac{ \pi_{\theta} (y_{i,t} | x, y_{i,<t}) }{ \pi_{\theta_\text{old}} (y_{i,t} | x,y_{i,<t})}$. These weights are uneven and can fall anywhere within $(0, 1+\varepsilon]$ when $\widehat{A}_{i} > 0$ or within $[1-\varepsilon, +\infty)$ for $\widehat{A}_{i} < 0$, making them significant and subject to compounding effects throughout the training process. Such non-uniformity can introduce instability and unpredictable behaviors. In comparison, GSPO treats all tokens within a response uniformly by assigning them equal weights, thereby removing the instability associated with GRPO’s variable weighting.

\subsection{GSPO-token: A Token-level Objective Variant}

In contexts such as multi-turn RL, it may be preferable to perform advantage adjustments at a more granular level than entire sequences. Accordingly, we propose a token-level adaptation of GSPO, referred to as \textbf{GSPO-token}, which enables the customization of advantages at the individual token level:

\begin{align}
\mathcal{J}_\text{GSPO-token}(\theta)
=
\mathbb{E}_{ x \sim \mathcal{D},\, \{y_i\}_{i=1}^G \sim \pi_{\theta_\text{old}}( \cdot | x) }
\left[ 
\frac{1}{G} \sum_{i=1}^{G} \frac{1}{|y_i|} \sum_{t=1}^{|y_i|} 
\min \left( s_{i,t}(\theta) \widehat{A}_{i,t},  \, \mathrm{clip} \left( s_{i,t}(\theta), 1 - {\varepsilon}, 1 + {\varepsilon} \right) \widehat{A}_{i,t} \right) 
\right],
\label{equ:gspo-token}
\end{align}

This objective, $\mathcal{J}_\text{GSPO-token}(\theta)$, is formulated as the expectation, over data samples $x$ drawn from $\mathcal{D}$ and generation batches $\{y_i\}_{i=1}^G$ sampled from the reference policy $\pi_{\theta_\text{old}}(\cdot|x)$, of the following inner sum: For each example $i$, a token-wise average is computed over its sequence of length $|y_i|$, and then the result is averaged across all $G$ examples in the batch. The summand at each timestep $t$ involves the minimum between the unaltered surrogate advantage term $s_{i,t}(\theta) \widehat{A}_{i,t}$ and its clipped counterpart, where $s_{i,t}(\theta)$ is confined to $[1-{\varepsilon}, 1+{\varepsilon}]$ via the $\mathrm{clip}$ operator.

where

\begin{align}
s_{i,t}(\theta) 
= \mathrm{sg} \left[ s_{i}(\theta) \right]  \cdot \frac{ \pi_{\theta} (y_{i,t} | x, y_{i,<t}) }{ \mathrm{sg} \left[ \pi_{\theta} (y_{i,t} | x, y_{i,<t}) \right] },
\end{align}

and $\mathrm{sg}[\cdot]$ denotes only taking the numerical value but stopping the gradient, corresponding to the \texttt{detach} operation in PyTorch. The gradient of GSPO-token can be derived as:

\begin{align}
\nabla_{\theta} \mathcal{J}_\text{GSPO-token}(\theta)
=&\ 
\nabla_{\theta} \mathbb{E}_{ x \sim \mathcal{D},\, \{y_i\}_{i=1}^G \sim \pi_{\theta_\text{old}}( \cdot | x) }
\left[ 
\frac{1}{G} \sum_{i=1}^{G}
\frac{1}{|y_i|} \sum_{t=1}^{|y_i|} 
s_{i,t}(\theta) \widehat{A}_{i,t}
\right] \\
=&\ 
\mathbb{E}_{ x \sim \mathcal{D},\, \{y_i\}_{i=1}^G \sim \pi_{\theta_\text{old}}( \cdot | x) }
\left[ 
\frac{1}{G} \sum_{i=1}^{G} s_i (\theta) \cdot \frac{1}{|y_i|} \sum_{t=1}^{|y_i|} 
\widehat{A}_{i,t} \frac{ \nabla_{\theta} \pi_{\theta} (y_{i,t} | x, y_{i,<t}) }{ \pi_{\theta} (y_{i,t} | x, y_{i,<t}) }
\right] \\
=&\ 
\mathbb{E}_{ x \sim \mathcal{D},\, \{y_i\}_{i=1}^G \sim \pi_{\theta_\text{old}}( \cdot | x) }
\left[ 
\frac{1}{G} \sum_{i=1}^{G}  \left( \frac{ \pi_{\theta} (y_i | x) }{ \pi_{\theta_\text{old}} (y_i | x)} \right)^{\frac{1}{|y_i|}}
\cdot \frac{1}{|y_i|} \sum_{t=1}^{|y_i|}
\widehat{A}_{i,t} \nabla_{\theta} \log \pi_{\theta} (y_{i,t} | x, y_{i,<t})  
\right].
\label{equ:gspo-token_grad}
\end{align}

Observe that the expression $\frac{ \pi_{\theta} (y_{i,t} | x, y_{i,<t}) }{ \mathrm{sg} \left[ \pi_{\theta} (y_{i,t} | x, y_{i,<t}) \right] }$ evaluates to 1, resulting in $s_{i,t}(\theta)$ being numerically equivalent to $s_{i}(\theta)$.
By examining Equation~\eqref{equ:gspo} alongside~\eqref{equ:gspo-token}, as well as Equation~\eqref{equ:gspo_grad} with~\eqref{equ:gspo-token_grad}, it follows that GSPO-token and GSPO yield identical results with respect to the optimization objective, clipping constraint, and the theoretical gradient, provided that the advantages assigned to all tokens in the sequence $y_i$ are held constant (specifically, $\widehat{A}_{i,t} = \widehat{A}_{i}$ for all $t$). However, GSPO-token enables a greater degree of granularity, allowing for the assignment of token-level advantage values.

\section{Experiments and Discussion}

\subsection{Empirical Results}

We conduct experiments using a cold-start model derived from Qwen3-30B-A3B-Base and present both the trajectories of training rewards and evaluation metrics across several benchmarks: AIME'24 (measured by average Pass@1 with 32 samples), LiveCodeBench (202410-202502, using average Pass@1 across 8 samples), and CodeForces (using Elo Rating). For reinforcement learning, the rollout data within each batch is divided into four mini-batches for performing gradient descent steps. Within the GSPO framework, the left and right clipping parameters in Equation~\eqref{equ:gspo} are set at 3e-4 and 4e-4, respectively. For comparison, we use GRPO as the baseline, adjusting its left and right clipping parameters in Equation~\eqref{equ:grpo} to 0.2 and 0.27—values finely calibrated for an equitable evaluation. It is important to point out that GRPO requires the Routing Replay approach to achieve stable MoE RL training, an aspect further elaborated in \S~\ref{subsec:routing_replay}. In contrast, \textit{GSPO does not depend on this methodology}.

As depicted in Figure~\ref{fig:results}, the training process utilizing GSPO remains stable at all stages. Notably, \textbf{GSPO facilitates ongoing improvements in model performance by leveraging increased training computation, systematically refreshing the query set, and generating longer outputs}. In addition, GSPO outperforms GRPO in terms of training efficiency, attaining higher training accuracy and superior benchmark results with equivalent computational resources and query budgets. Lastly, GSPO has been effectively implemented in the RL training regimen for the latest Qwen3 models, further substantiating its effectiveness in harnessing RL scaling capabilities for large language models.

\subsection{Curious Observation on Clipping Fractions}

GSPO is notably different from GRPO in that it applies clipping to whole responses instead of at the level of individual tokens. As illustrated in Figure~\ref{fig:clipping}, this strategy results in GSPO clipping two orders of magnitude more tokens than GRPO, and this pronounced gap persists regardless of the specific clipping ranges used. Interestingly, even though GSPO discards a much larger proportion of tokens—thus having less data available for training or gradient estimation—it nonetheless demonstrates higher training efficiency than GRPO. This somewhat surprising outcome—that greater token-level clipping can lead to improved training outcomes—suggests that GRPO's approach of estimating gradients at the token level is fundamentally noisier and less effective in utilizing samples. By focusing on sequence-level clipping, GSPO is able to generate a more dependable and informative learning signal.

\subsection{Benefit of GSPO for MoE Training}
\label{subsec:routing_replay}

\paragraph{Background}
In contrast to RL training with dense architectures, the inherent sparsity in MoE models leads to distinct issues affecting training stability. Notably, our investigations revealed that leveraging the GRPO algorithm exposes MoE models to \textit{expert-activation volatility}, which may obstruct proper RL convergence. Specifically, following one or more steps of gradient optimization, the ensemble of experts selected for identical responses can differ markedly. For instance, when using the 48-layer Qwen3-30B-A3B-Base model, every RL gradient step yields about a 10\% turnover in active experts for the same rollout instance when comparing the new policy $\pi_{\theta}$ against the former policy $\pi_{\theta_\text{old}}$.
Such volatility is particularly pronounced as model depth increases, which in turn induces dramatic instability in the token-level importance ratios, $w_{i,t}(\theta) = \frac{ \pi_{\theta} (y_{i,t} | x, y_{i,<t}) }{ \pi_{\theta_\text{old}} (y_{i,t} | x,y_{i,<t})}$. As outlined in \S~\ref{sec:motivation} and~\ref{subsec:gradient}, this undermines the reliability of these ratios, thereby impeding the typical convergence of RL training routines.

\paragraph{Our Previous Approach} In addressing this issue, our earlier work introduced the \textbf{Routing Replay} training method. This involved saving the set of experts activated under $\pi_{\theta_\text{old}}$ and subsequently “replaying” these routing decisions within $\pi_{\theta}$ when evaluating the importance weights $w_{i,t}(\theta) = \frac{ \pi_{\theta} (y_{i,t} | x, y_{i,<t}) }{ \pi_{\theta_\text{old}} (y_{i,t} | x,y_{i,<t})}$. By ensuring that, for each token $y_{i,t}$, both $\pi_{\theta} (y_{i,t} | x, y_{i,<t})$ and $\pi_{\theta_\text{old}} (y_{i,t} | x,y_{i,<t})$ make use of an identical routing structure, we are able to maintain stability in the computed token-level importance ratios, thereby facilitating optimization over the same routed sub-network through different gradient steps. Figure~\ref{fig:routing_replay} illustrates how Routing Replay functions as a pivotal component for ensuring stable convergence in the GRPO-based training of MoE architectures.

\paragraph{Benefit of GSPO} While Routing Replay makes it possible for MoE models trained with GRPO to reach convergence, its strategy of duplicating routing configurations adds extra burdens in terms of memory usage and communication, and it restricts the MoE model from utilizing its true representational capacity. In comparison, as illustrated in Figure~\ref{fig:results}, GSPO completely avoids relying on Routing Replay. Instead, it is able to directly compute the importance ratios $s_i(\theta)$ through standard means, achieving steady convergence and consistent optimization. The central observation is that GSPO relies exclusively on the sequence-level likelihood (i.e., $\pi_{\theta} (y_i | x)$), rendering it largely indifferent to the token-level likelihood (i.e., $\pi_{\theta} (y_{i,t} | x, y_{i,<t})$). Because the fundamental language modeling proficiency of the MoE model remains intact, fluctuations in sequence likelihood are minimal. In essence, GSPO provides a principled solution to the instability of expert activation in MoE models, thereby eliminating the necessity for complex techniques such as Routing Replay. This approach not only streamlines and strengthens the training regimen but also unlocks the full capacity of the model by removing imposed constraints.

\subsection{Benefit of GSPO for RL Infrastructure}

Due to differences in numerical precision between training platforms (such as Megatron) and inference systems (like SGLang and vLLM), it is common in real-world applications to recalculate the likelihoods of generated responses under the previous policy $\pi_{\theta_\text{old}}$ using the training engine. In contrast, GSPO relies on sequence-level rather than token-level likelihoods for its optimization procedure; sequence-level metrics are inherently less sensitive to minor discrepancies in precision. As a result, GSPO enables the direct utilization of likelihood values produced by the inference engine during optimization, eliminating the requirement to recompute these with the training engine. This feature is particularly advantageous in contexts involving partial rollouts, multi-turn reinforcement learning, and frameworks that separate training and inference processes.

\section{Conclusion}

Group Sequence Policy Optimization (GSPO) is introduced as a novel reinforcement learning algorithm designed for the training of large language models. GSPO utilizes importance sampling principles by establishing importance ratios grounded in sequence likelihood and employs sequence-level strategies for clipping, reward assignment, and optimization. Relative to GRPO, GSPO offers substantially improved training stability, computational efficiency, and overall performance. It is particularly well-suited for large-scale reinforcement learning tasks involving MoE models, thereby playing a pivotal role in the advancements realized in recent Qwen3 models. With GSPO serving as a scalable foundational technique, further scaling of reinforcement learning is anticipated to drive significant progress in artificial intelligence.

\end{document}